Transit accessibility is among the sharpest markers of community type in the United States: transit-rich urban census tracts share infrastructure, commuting behaviour, and policy interests fundamentally different from car-dependent suburban or rural tracts. This paper operationalises transit-accessibility similarity as an edge weight modifier for the \textsc{bisect} redistricting pipeline. Using General Transit Feed Specification (GTFS) data from regional transit agencies, we compute a per-tract transit score ($\sum_{\text{routes within 400m}} \text{frequency} \times \text{speed factor}$) and derive a similarity metric bounded in $[0,1]$. Adjacent tracts with similar transit access receive heavy edge weights; transit-rich urban tracts adjacent to car-dependent suburban tracts receive light weights, inviting METIS to cut between them. Applied to North Carolina, Wisconsin, and Texas, transit-accessibility weights keep urban transit corridors together (mean within-transit-zone edge weight increase: $+31\%$) while allowing cuts at transit transition boundaries (urban-suburban edges reduced to $62\%$ of baseline weight). Rural-to-rural pairs where both tracts have zero transit score maintain baseline weight, preventing the signal from disadvantaging rural communities. We implement transit weights as the \texttt{--weights-override transit} flag in the \textsc{bisect} CLI, using only publicly available GTFS data from transit agencies' official feeds. Legal grounding is the communities-of-interest doctrine applied to shared infrastructure and commuting patterns as community identity markers.
